\section{Limitations}
\label{sec:limitations}

\paragraph{Crowded scenes}
Synthetic pre-training spans one to five people.
COCO fine-tuning and the K-head do encounter larger scenes, but such scenes form a minority:
12.5\,\% of the \texttt{val2017} images with keypoint annotations contain six or more people.
On those images the K-head predicts the exact count in 22.9\,\% of scenes
and under-counts in 69.4\,\%, and the AP gap to the native grouping widens from 0.029 to 0.076.
Extending the person-count range of the synthetic data is one direct remedy.

\paragraph{Pose grouping accuracy and system accuracy}
The primary metric scores the grouping of matched detections,
and in that protocol the graph handed to PG-GAT is likewise built from matched detections;
false-positive keypoints therefore never reach it.
The system-level AP evaluation lifts both restrictions
and shows a larger cost than the matched-detection accuracy suggests.
PGA remains the right instrument for locating the bottleneck in the grouping stage,
but it is not an estimate of deployed performance.

\paragraph{Extremities and the skeleton prior}
On multi-person scenes the misassignment rate of PG-GAT with the annotated count rises from 5--6\,\% at the shoulders to 10.6--12.1\,\% at the ankles,
while the native grouping stays within 2.6\,\% to 4.4\,\% for every joint type.
The skeleton that supplies the edge categories and hop distances omits the links between the head and the shoulders;
head joints are nevertheless grouped more accurately than body joints,
and the omission therefore does not explain the deficit.
Structure beyond individual joints, such as limbs, is not represented in the model.

\paragraph{Single runs}
The ablation, the sweeps, and the variants are single training runs,
with no repeated seeds or confidence intervals.
Differences of the size of the individual modification gains, or of the 0.0020 fine-tuning-sweep improvement,
are not established as systematic.
The graph's neighbourhood size was left out of every sweep.

\paragraph{Detector comparisons}
The matched subsets differ in size and recall from one detector to another,
and accuracies are comparable within a detector but only with caution across detectors.
The RTMO-L native grouping draws on weights trained over a multi-dataset corpus.
The AP evaluation runs a single-scale HigherHRNet wrapper without flip testing,
and its absolute values fall below the detector's published figure.